\textbf{Phần 1. Trắc nghiệm lựa chọn}
\setcounterpageref{ex}{0}

\begin{ex} %[4D10V56]
Rút gọn phân số là
\choice
{Nhân cả tử và mẫu với một số}
{Cộng thêm vào tử và mẫu}
{Đổi chỗ tử số và mẫu số}
{\True Chia cả tử số và mẫu số cho cùng một số tự nhiên lớn hơn 1 để được phân số đơn giản hơn và bằng phân số đã cho}
\loigiai{
Rút gọn phân số là chia cả tử và mẫu cho ước chung lớn hơn 1.\\
\textbf{Đáp án của câu hỏi là D}
}
\end{ex}

\begin{ex} %[4D10V56]
Phân số tối giản là phân số
\choice
{Có tử số bằng 1}
{\True Tử số và mẫu số không cùng chia hết cho một số tự nhiên nào lớn hơn 1}
{Có mẫu số bằng 10}
{Có tử số nhỏ hơn 10}
\loigiai{
Phân số tối giản không thể rút gọn được nữa.\\
\textbf{Đáp án của câu hỏi là B}
}
\end{ex}

\begin{ex} %[4D10V56]
Rút gọn phân số $\dfrac{6}{8}$ về dạng tối giản ta được
\choice
{$\dfrac{2}{4}$}
{$\dfrac{3}{2}$}
{\True $\dfrac{3}{4}$}
{$\dfrac{1}{2}$}
\loigiai{
$\dfrac{6 : 2}{8 : 2} = \dfrac{3}{4}$.\\
\textbf{Đáp án của câu hỏi là C}
}
\end{ex}

\begin{ex} %[4D10V56]
Trong các phân số sau, phân số nào là phân số tối giản?
\choice
{\True $\dfrac{5}{7}$}
{$\dfrac{4}{6}$}
{$\dfrac{9}{12}$}
{$\dfrac{15}{20}$}
\loigiai{
$\dfrac{5}{7}$ là phân số tối giản.\\
\textbf{Đáp án của câu hỏi là A}
}
\end{ex}

\begin{ex} %[4D10V56]
Rút gọn phân số $\dfrac{18}{27}$ ta được phân số tối giản là
\choice
{$\dfrac{6}{9}$}
{$\dfrac{3}{4}$}
{$\dfrac{1}{3}$}
{\True $\dfrac{2}{3}$}
\loigiai{
$\dfrac{18 : 9}{27 : 9} = \dfrac{2}{3}$.\\
\textbf{Đáp án của câu hỏi là D}
}
\end{ex}

\begin{ex} %[4D10V56]
Phân số nào sau đây CHƯA tối giản?
\choice
{$\dfrac{3}{5}$}
{\True $\dfrac{14}{21}$}
{$\dfrac{7}{9}$}
{$\dfrac{11}{13}$}
\loigiai{
$\dfrac{14}{21}$ cùng chia hết cho 7 (rút gọn được thành $\dfrac{2}{3}$).\\
\textbf{Đáp án của câu hỏi là B}
}
\end{ex}

\begin{ex} %[4D10V56]
Rút gọn phân số $\dfrac{45}{60}$ về tối giản ta được
\choice
{$\dfrac{9}{12}$}
{$\dfrac{15}{20}$}
{\True $\dfrac{3}{4}$}
{$\dfrac{4}{5}$}
\loigiai{
$\dfrac{45 : 15}{60 : 15} = \dfrac{3}{4}$.\\
\textbf{Đáp án của câu hỏi là C}
}
\end{ex}

\begin{ex} %[4D10V56]
Rút gọn phân số $\dfrac{75}{100}$ ta được
\choice
{\True $\dfrac{3}{4}$}
{$\dfrac{15}{20}$}
{$\dfrac{7}{10}$}
{$\dfrac{4}{5}$}
\loigiai{
$\dfrac{75 : 25}{100 : 25} = \dfrac{3}{4}$.\\
\textbf{Đáp án của câu hỏi là A}
}
\end{ex}

\begin{ex} %[4D10V56]
Cho phân số $\dfrac{36}{48}$. Số tự nhiên lớn nhất có thể dùng để rút gọn ngay phân số này về tối giản là
\choice
{$6$}
{$4$}
{$8$}
{\True $12$}
\loigiai{
Ước chung lớn nhất của 36 và 48 là 12.\\
\textbf{Đáp án của câu hỏi là D}
}
\end{ex}

\begin{ex} %[4D10V56]
Trong các phân số $\dfrac{12}{15}, \dfrac{7}{9}, \dfrac{20}{25}, \dfrac{11}{14}$, có bao nhiêu phân số tối giản?
\choice
{$1$ phân số}
{\True $2$ phân số}
{$3$ phân số}
{$4$ phân số}
\loigiai{
Các phân số tối giản là: $\dfrac{7}{9}$ và $\dfrac{11}{14}$ (có 2 phân số).\\
\textbf{Đáp án của câu hỏi là B}
}
\end{ex}

\begin{ex} %[4D10V56]
Rút gọn phân số $\dfrac{120}{150}$ về dạng tối giản:
\choice
{$\dfrac{12}{15}$}
{$\dfrac{6}{7}$}
{\True $\dfrac{4}{5}$}
{$\dfrac{2}{3}$}
\loigiai{
$\dfrac{120 : 30}{150 : 30} = \dfrac{4}{5}$.\\
\textbf{Đáp án của câu hỏi là C}
}
\end{ex}

\begin{ex} %[4D10V56]
Tính giá trị biểu thức sau bằng cách rút gọn: $\dfrac{2 \times 3 \times 5 \times 7}{3 \times 5 \times 7 \times 8}$.
\choice
{\True $\dfrac{1}{4}$}
{$\dfrac{1}{2}$}
{$\dfrac{2}{8}$}
{$1$}
\loigiai{
Cùng chia tử và mẫu cho $3 \times 5 \times 7$: được $\dfrac{2}{8} = \dfrac{1}{4}$.\\
\textbf{Đáp án của câu hỏi là A}
}
\end{ex}

\textbf{Phần 2. Trắc nghiệm đúng sai}
\setcounterpageref{ex}{0}

\begin{ex} %[4D10V56]
Xét tính đúng sai về rút gọn phân số:
\choiceTF
{\True $\dfrac{8}{12}$ rút gọn được thành $\dfrac{2}{3}$.}
{\True $\dfrac{13}{17}$ là phân số tối giản.}
{Phân số $\dfrac{15}{35}$ là phân số tối giản.}
{\True $\dfrac{15}{35}$ rút gọn được thành $\dfrac{3}{7}$.}
\loigiai{
\begin{itemize}
    \item $8/12 = 2/3$. Đúng.
    \item $13/17$ tối giản: Đúng.
    \item $15/35$ chia hết cho 5 nên chưa tối giản. Sai.
    \item $15/35 = 3/7$. Đúng.
\end{itemize}
\begin{itemize}
    \item a) \textbf{Đúng}.
    \item b) \textbf{Đúng}.
    \item c) \textbf{Sai}.
    \item d) \textbf{Đúng}.
\end{itemize}
}
\end{ex}

\begin{ex} %[4D10V56]
Cho phân số $\dfrac{40}{60}$. Xét tính đúng sai:
\choiceTF
{\True Chia cả tử và mẫu cho 10 được $\dfrac{4}{6}$.}
{\True Rút gọn về phân số tối giản được $\dfrac{2}{3}$.}
{Rút gọn về phân số tối giản được $\dfrac{4}{6}$.}
{\True Phân số $\dfrac{40}{60}$ bằng phân số $\dfrac{10}{15}$.}
\loigiai{
\begin{itemize}
    \item Chia cho 10 được 4/6: Đúng.
    \item Tối giản là 2/3: Đúng.
    \item 4/6 chưa phải tối giản. Sai.
    \item $40/60 = 10/15 = 2/3$. Đúng.
\end{itemize}
\begin{itemize}
    \item a) \textbf{Đúng}.
    \item b) \textbf{Đúng}.
    \item c) \textbf{Sai}.
    \item d) \textbf{Đúng}.
\end{itemize}
}
\end{ex}

\textbf{Phần 3. Trả lời ngắn}
\setcounterpageref{ex}{0}

\begin{ex} %[4D10V56]
Rút gọn phân số $\dfrac{10}{15}$ về tối giản (dạng a/b).
\shortans{2/3}
\loigiai{
\textbf{Đáp số:} $2/3$\\
$2/3$.
}
\end{ex}

\begin{ex} %[4D10V56]
Rút gọn phân số $\dfrac{24}{32}$ về tối giản (dạng a/b).
\shortans{3/4}
\loigiai{
\textbf{Đáp số:} $3/4$\\
$3/4$.
}
\end{ex}

\begin{ex} %[4D10V56]
Rút gọn phân số $\dfrac{42}{56}$ về tối giản (dạng a/b).
\shortans{3/4}
\loigiai{
\textbf{Đáp số:} $3/4$\\
$42 : 14 / 56 : 14 = 3/4$.
}
\end{ex}

\begin{ex} %[4D10V56]
Rút gọn phân số $\dfrac{50}{125}$ về tối giản (dạng a/b).
\shortans{2/5}
\loigiai{
\textbf{Đáp số:} $2/5$\\
$50 : 25 / 125 : 25 = 2/5$.
}
\end{ex}

\begin{ex} %[4D10V56]
Tính giá trị rút gọn của phân số $\dfrac{16 \times 5}{25 \times 4}$ (dạng a/b).
\shortans{4/5}
\loigiai{
\textbf{Đáp số:} $4/5$\\
$(16 : 4) / (25 : 5) = 4/5$.
}
\end{ex}

\begin{ex} %[4D10V56]
Rút gọn biểu thức: $\dfrac{12 \times 15 \times 18}{24 \times 30 \times 36}$ về phân số tối giản (dạng a/b).
\shortans{1/8}
\loigiai{
\textbf{Đáp số:} $1/8$\\
$(1/2) \times (1/2) \times (1/2) = 1/8$.
}
\end{ex}
